\documentclass[11pt]{article}

% --- Packages (all standard; compiles with a base TeX Live install) --------
\usepackage[margin=1in]{geometry}
\usepackage{amsmath,amssymb,amsthm}
\usepackage{booktabs}
\usepackage{array}
\usepackage{enumitem}
\usepackage[hidelinks]{hyperref}
\usepackage{titlesec}

\titleformat{\section}{\normalfont\large\bfseries}{\thesection}{0.6em}{}
\titleformat{\subsection}{\normalfont\normalsize\bfseries}{\thesubsection}{0.6em}{}

% Lightweight proposition environment (no fake theorems — claims with proofs
% or clear justification only).
\newtheorem{proposition}{Proposition}
\newtheorem{observation}{Observation}

\title{\textbf{Project Lifeline}\\[3pt]
\large Offline, End-to-End-Encrypted Emergency Messaging,\\
and What Composing It Taught Us About Metering a Partitioned Network}
\author{Archit Sharma\thanks{Project Lifeline. Correspondence: \texttt{archit.sharma@nometria.com}.
Source: \url{https://github.com/matrix-share/matrix}.}}
\date{\today}

\begin{document}
\maketitle

\begin{abstract}
Disasters remove the infrastructure messaging depends on: towers, backhaul,
power, and the reachable servers that hold our identities and authorize our
traffic. \emph{Project Lifeline} is an implemented, decentralized, offline-first,
end-to-end-encrypted messenger built as a Delay/Disruption-Tolerant Network
(DTN)~\cite{rfc4838}. We are candid up front about what it is and is not:
\emph{the system invents no cryptographic primitive and few mechanisms} --- it
composes self-certifying identity, content-addressed transfer, spray-and-wait
routing, sealed-sender encryption, and macaroon-style capabilities, all
off-the-shelf. Its research interest is not the parts but \emph{three
observations that emerged from putting them together at the DTN$\times$security
boundary}. The central one: metering a scarce shared resource --- internet egress
through a gateway --- with an offline-verifiable bearer capability in a
partitioned network is an instance of \emph{offline double-spending}. This bounds
what any such system can guarantee (a global budget is unenforceable without
consensus; only a per-partition-component bound is achievable) and points to a
Chaum--Fiat--Naor-style design that \emph{detects and attributes} over-spend when
partitions merge rather than pretending to prevent it. The second: under
sealed-sender delivery, reputation \emph{must} be source-local, which
incidentally makes it immune to the Sybil badmouthing that plagues shared
reputation --- so confidentiality and accountability align rather than conflict.
The third is an architectural principle --- an open, un-authorized relay commons
beneath a gated, capability-metered egress. We describe the system, state these
three observations as propositions, give a threat-model security analysis and a
reproducible comparative evaluation, and are explicit throughout about which
guarantees are proven, which are engineering, and which remain open (chief among
them an independent cryptographic audit).
\end{abstract}

\tableofcontents

% ===========================================================================
\section{Introduction}
% ===========================================================================

The design center of most messaging systems is a reachable server. In a
disaster --- an earthquake, a flood, a grid failure, a deliberate shutdown ---
that assumption is exactly what fails. Project Lifeline targets the regime where
connectivity is \emph{intermittent, partitioned, and heterogeneous}: some nodes
share a short-range radio for minutes, some physically carry data between
partitions, and a lucky few briefly reach the wider internet.

Following the end-to-end argument~\cite{saltzer1984endtoend} and the DTN
architecture~\cite{rfc4838,rfc5050,rfc9171}, our stance is that reliability,
security, and \emph{authorization} must live at the endpoints and in the data
they carry, never in the links or in a server one must reach. The bulk of the
paper (\S\S\ref{sec:arch}--\ref{sec:egress}) shows how each subsystem realizes
that stance by composing known techniques; we flag honestly where each part is
standard.

\paragraph{What is actually new.} The contribution is \emph{not} a new primitive
or protocol. It is three observations that only became visible once the pieces
were assembled, developed in \S\ref{sec:obs}:
\begin{enumerate}[leftmargin=1.6em,itemsep=1pt]
  \item \textbf{Egress metering is offline double-spending}
  (\S\ref{sub:doublespend}, Proposition~\ref{prop:doublespend}). Offline-verifiable
  bearer capabilities cannot enforce a \emph{global} usage budget across a
  partition; the achievable bound is per-component, and the principled response
  is detection-and-attribution on merge, not prevention.
  \item \textbf{Sealed-sender aligns privacy with accountability}
  (\S\ref{sub:privacy}, Proposition~\ref{prop:privacy}). Hiding the sender forces
  reputation to be source-local, which removes the shared state a Sybil would
  badmouth.
  \item \textbf{Open relay, gated egress} (\S\ref{sub:asymmetry},
  Observation~\ref{obs:asym}) --- the asymmetry that organizes the whole design
  and from which the first observation follows.
\end{enumerate}
Everything else in the paper is engineering in service of these, plus an honest
account (\S\ref{sub:composition}) of what is off-the-shelf.

\paragraph{Status.} Lifeline is an implemented system: a Rust workspace of 15
crates, $\sim$210 automated tests, and an acceptance simulation held to a
$\geq 95\%$ delivery target. It has \emph{not} had an independent cryptographic
review; where a guarantee is aspirational we say so.

% ===========================================================================
\section{System and Threat Model}\label{sec:model}
% ===========================================================================

\paragraph{Network.} The network is a time-varying graph whose edges (a BLE
contact, Wi-Fi Aware, a LoRa hop, an ultrasonic burst, an internet tunnel, a
bridged relay) appear and vanish. There is no assumption of a contemporaneous
end-to-end path; delivery is store-carry-forward. Each edge type is an
\textsc{Interface} with an MTU and a throughput class. The network may be
\emph{partitioned} into components that cannot communicate for extended periods.

\paragraph{Adversary.} A Dolev--Yao network adversary can observe, drop, replay,
reorder, and inject on any link. A fraction of \emph{participants} are malicious
but protocol-conformant: they may drop bundles they accept, forge control
messages, mint many identities (Sybil), or spam. Endpoint devices and long-term
keys are assumed uncompromised; seizure/coercion caveats are discussed in
\S\ref{sec:security}. We do \emph{not} claim resistance to a global traffic
analyst~\cite{vandenhooff2015vuvuzela,piotrowska2017loopix}.

\paragraph{Goals, in priority order.} (1)~life-safety delivery --- an SOS moves
even under spam; (2)~confidentiality and integrity of content and, as far as
possible, metadata; (3)~availability under partition; (4)~augmentation, not
replacement, of existing networks.

% ===========================================================================
\section{Architecture}\label{sec:arch}
% ===========================================================================

Lifeline is five layers, each a thin trait boundary (Table~\ref{tab:layers}).
The invariant that makes it tractable: \emph{relays are dumb}. An intermediate
node forwards an opaque, sealed \textsc{Bundle}; it can neither read it nor is
trusted to have handled it correctly. All correctness is re-established at the
recipient. The narrow waist is the sealed bundle (cf.\ the IP hourglass); the
engine drives every attached \textsc{Interface} concurrently and deduplicates by
bundle id, so a message riding BLE \emph{and} LoRa is stored once. A new carrier
or bridged network is one trait implementation --- proven by two live external
bridges (Nostr, Meshtastic, \S\ref{sec:bridge}).

\begin{table}[t]\centering\small
\begin{tabular}{@{}p{2.2cm}p{4.2cm}p{6.4cm}@{}}
\toprule
\textbf{Layer} & \textbf{Concern} & \textbf{Key abstraction / crate}\\
\midrule
L1 Identity & self-certifying naming & \textsc{Address}$=\mathrm{BLAKE3}(\mathrm{pk})[{:}16]$ (\texttt{proto}, \texttt{core})\\
L2 Bundle & sealed unit of transfer & \textsc{Bundle}, \textsc{Manifest}, content id (\texttt{proto}, \texttt{core})\\
L3 Discovery & find peers / content & Kademlia DHT, geocast, time-sync (\texttt{dht}, \texttt{geo}, \texttt{timesync})\\
L4 Routing & store-carry-forward & spray-and-wait, custody, attribution (\texttt{router}, \texttt{engine})\\
L5 Application & messages, groups, egress & sealed-sender, groups, gateway (\texttt{core}, \texttt{inet})\\
\bottomrule
\end{tabular}
\caption{The five layers. The \textsc{Interface}, \textsc{ExternalNet}, and
\textsc{AccessPolicy} traits are the extension seams: a new carrier, bridged
network, or egress policy is one \texttt{impl}.}
\label{tab:layers}
\end{table}

% ===========================================================================
\section{Identity and the Bundle (standard)}\label{sec:bundle}
% ===========================================================================

\paragraph{Self-certifying identity.} A node generates an Ed25519 signing key and
an X25519 key-agreement key on first launch. Its \textsc{Address} is
$\mathrm{BLAKE3}(\text{sign\_pub})[{:}16]$ --- a 128-bit fingerprint of the public
key. There is no registrar: the key \emph{is} the identity, which is what lets
identity survive when servers do not. Addresses double as the routing destination
and the Kademlia key. This is standard self-certifying naming.

\paragraph{Content addressing and the manifest.} Large objects are split into
fixed-size blocks, each named by its BLAKE3 hash --- its content id
(CID)~\cite{benet2014ipfs}. An ordered \textsc{Manifest} (named by a root CID over
length, block size, and block CIDs) is a Merkle description of the whole
object~\cite{merkle1987digsig}. This buys, for free: \emph{intrinsic integrity}
(a block is trusted only if it hashes to the requested CID), \emph{global
deduplication} (identical content $\Rightarrow$ identical CID $\Rightarrow$ stored
once, over any number of carriers), and a \emph{completeness predicate} (the
object is done iff the set of missing manifest CIDs is empty). For cross-partition
resilience a bundle may additionally be Reed--Solomon erasure
coded~\cite{reed1960polynomial}. All standard.

% ===========================================================================
\section{Routing (standard, with one twist)}\label{sec:routing}
% ===========================================================================

\paragraph{Store-carry-forward with binary spray-and-wait.} Routing is binary
spray-and-wait~\cite{spyropoulos2005sprayandwait}: a bundle carries a copy budget
$L$; a node holding $c>1$ copies gives half to a new contact and keeps half; a
node holding one copy waits to deliver directly. This recovers most of epidemic
routing's~\cite{vahdat2000epidemic} delivery probability at a fraction of its
transmission overhead (quantified in \S\ref{sec:eval}). Standard.

\paragraph{Priority and anti-spam postage.} Bundles carry a priority
$\in\{$SOS, Alert, Normal, Bulk$\}$. Normal/Bulk bundles must carry Hashcash-style
proof-of-work postage~\cite{back2002hashcash} whose difficulty scales inversely
with priority; SOS is exempt. A relay verifies postage before admitting a
low-priority bundle. Standard.

\paragraph{The twist: source-attributed custody reputation.}\label{par:custody}
Custody transfer~\cite{rfc5050} lets a committed custodian accept a bundle so an
upstream node frees its copy --- creating a \emph{black-hole} risk (accept
custody, silently drop). Detecting this in a DTN is hard: no path, no timely ACK.
Our mechanism exploits a structural fact of sealed-sender delivery: a delivery
receipt is decryptable \emph{only by the original sender}, so \emph{only the
source can attribute}. The source ledgers $\text{bundle}\!\rightarrow\!
\text{custodian}$ at each hand-off; a verified receipt credits that bundle's
custodians; a custodian whose carried bundles repeatedly fail to deliver --- both
by a consecutive-miss count (a black hole) and by a lifetime delivery ratio (a
grey hole that drops selectively) --- is demoted. Demotion only \emph{re-routes
around} a peer when an alternative exists; it never drops traffic, so the
$\geq 95\%$ target is safe even when the heuristic is wrong. This source-locality
is not merely convenient; it is forced by confidentiality, and it has a security
consequence we return to in \S\ref{sub:privacy}. Implemented in
\texttt{router::attribution}.

% ===========================================================================
\section{Reliable Multi-Carrier Transfer (standard)}\label{sec:swarm}
% ===========================================================================

The disaster requirement --- an object rides every carrier at once and the
recipient ends up with all of it, exactly once, over lossy intermittent links ---
falls out of the pieces already in place: content addressing gives
no-duplication and integrity, the manifest gives a completeness predicate, and
the multi-interface engine drives every carrier. The one gap is \emph{multi-source
acquisition}, closed by a manifest-scoped HAVE exchange adapting the BitTorrent
swarm~\cite{cohen2003bittorrent} to a DTN: the fetcher learns which provider holds
which missing CID (a per-object HAVE bitmap), requests each block from a holder,
and rotates a still-missing block to a \emph{different} provider on retry ---
routing around a dark or black-hole provider. At most one provider is asked per
block per round (no duplicate traffic); content addressing makes duplicate
arrivals idempotent; the missing set is monotone, so completion is provable. For
\emph{whole-store} anti-entropy between peers we instead exchange a fingerprinted
held-bundle digest, retiring redundant offers, with range-based set
reconciliation~\cite{minsky2003setrec,eppstein2011difference} as the
larger-scale successor. All standard building blocks.

% ===========================================================================
\section{Cryptographic Design (standard, with an honest caveat)}\label{sec:crypto}
% ===========================================================================

Lifeline composes audited primitives --- Ed25519, X25519, HKDF-SHA256,
XChaCha20-Poly1305, BLAKE3~\cite{oconnor2020blake3}, Argon2id~\cite{biryukov2016argon2}
--- and invents none.

\paragraph{Sealed-sender envelope.} A message is an ephemeral-static sealed box
to the recipient with the \emph{sender's signature inside the sealed plaintext},
so relays learn neither content nor sender. Forward secrecy comes not from a live
ratchet (impractical when a sender is offline for days) but from a ring of
rotating prekeys that bounds the window a seized long-term key can decrypt. This
is a deliberate, \emph{documented weakening} relative to a full Double
Ratchet~\cite{perrin2016doubleratchet}; it approximates X3DH-style asynchronous
agreement~\cite{marlinspike2016x3dh} for the store-and-forward setting. Groups use
sender-keys; an optional onion layer~\cite{dingledine2004tor} hides the
communication graph from any single relay.

% ===========================================================================
\section{Discovery and Bridges (standard)}\label{sec:discovery}
% ===========================================================================

When some nodes have connectivity, a Kademlia DHT~\cite{maymounkov2002kademlia}
over the 128-bit address space locates mailboxes and providers; it is
transport-agnostic and explicitly not relied on offline.
Geocast~\cite{niemeyer2008geohash} addresses ``everyone in cell $g$'' via a
region-derived key --- an addressing convenience, not a privacy mechanism (the
region commitment is reversible, stated candidly). With no reachable
NTP~\cite{rfc5905}, nodes acquire a coarse ``mesh time'' from stratum-tagged
beacons with bounded-slew correction, disciplined \emph{only} by explicitly
paired contacts and used advisorily (never for TTL, since beacons are
unauthenticated).

\label{sec:bridge}Lifeline augments rather than replaces: the
\textsc{ExternalNet} trait makes any external network look like another carrier.
Two bridges are implemented --- a live Nostr relay client~\cite{nostrnip01} (with
an unlinkable, domain-separated bearer subkey) and a Meshtastic MQTT
adapter~\cite{meshtastic} --- both off by default.

% ===========================================================================
\section{Capability-Governed Internet Egress}\label{sec:egress}
% ===========================================================================

A node with real internet can be an \emph{exit} for others: the mesh becomes the
access network, that node the gateway, over a store-and-forward request/response
proxy (a live socket is impractical over a high-latency DTN). The governing
requirement is a split we return to in \S\ref{sub:asymmetry}: \emph{mesh messages
are open to all; internet egress is node-authorized}.

An allow-list on the gateway fails under partition: the gateway cannot reach
whatever server provisions the list. Every comparable real system keeps the
authority in a reachable server --- cellular PCRF, captive-portal RADIUS,
enterprise zero-trust proxies~\cite{rose2020ztarch}. Instead we put the authority
\emph{in a token the requester presents}, following
macaroons~\cite{birgisson2014macaroons} and Biscuit~\cite{biscuit2021} on the
Ed25519 we already ship. A capability is a signed grant plus a signature-chain of
narrowing attenuations: (i)~\emph{offline-verifiable} (a gateway trusting the
issuer checks a signature and a local clock --- no server); (ii)~\emph{narrow-only
delegable} (a holder appends a scope caveat and passes it on without contacting
the issuer; the verifier conjoins every level, so a caveat can only remove
authority); (iii)~\emph{unstrippable} (each attenuation is signed by the prior
level's delegated key and binds the prior signature). A \textsc{ServiceClass}
label ($\textsf{LifeSafety}\!\subset\!\textsf{Messaging}\!\subset\!\textsf{Interactive}
\!\subset\!\textsf{Bulk}$) tiers egress, the exit-side analogue of a network
slice / DiffServ class~\cite{rfc2475,foukas2017slicing}. The gateway enforces SSRF
defenses~\cite{owasp_ssrf}: \texttt{http}/\texttt{https} only; never loopback,
private, link-local, CGNAT~\cite{rfc6598}, or cloud-metadata; the real fetcher
re-checks the \emph{resolved} IP against DNS rebinding. Real-time calling is
blocked \emph{structurally} --- a request/response DTN has no representation for a
persistent media stream, so no deep packet inspection is needed.

Applying capabilities to egress is engineering. What makes it interesting is what
happens to their \emph{quota} under partition, which is the subject of the next
section.

% ===========================================================================
\section{What Composing This Taught Us}\label{sec:obs}
% ===========================================================================

This section states the paper's actual contribution: three observations that the
composition made visible, none of which is a new primitive.

\subsection{Open relay, gated egress}\label{sub:asymmetry}

\begin{observation}\label{obs:asym}
A disaster mesh has two resource classes with opposite economics.
\emph{Relaying} a sealed message is an un-metered commons: bandwidth spent
carrying someone's SOS is cheap, refusing to carry it violates the primary goal,
and a relay cannot even read what it carries --- so relay must be \emph{open and
un-authorized}. \emph{Internet egress} is the opposite: a scarce, node-owned,
externally-accountable resource (the operator pays for and is liable for it) ---
so egress must be \emph{gated}. The gate cannot depend on a reachable authority
(\S\ref{sec:egress}), so it must be an offline-verifiable capability.
\end{observation}

This asymmetry --- open L4 relay beneath gated L5 egress --- is the organizing
principle of the design. It is also what raises the question the rest of the
section answers: once the gate is an offline bearer token, \emph{what can its
usage budget guarantee?}

\subsection{Metering an offline capability is double-spending}\label{sub:doublespend}

Give a capability a global usage budget $B$ (say, bytes of egress) and let any
gateway verify and meter it offline. What can be guaranteed about total
consumption?

\begin{proposition}[Egress quota is offline double-spending]\label{prop:doublespend}
Let a bearer capability authorize a global budget $B$, verifiable and meterable
offline by any gateway that trusts the issuer, with no communication between
gateways at spend time. If the network partitions into components
$C_1,\dots,C_k$, a holder present in reach of a gateway in each component can
consume up to $B$ in \emph{each}, for a total up to $kB$. No purely offline scheme
can do better: enforcing a global bound below $kB$ requires two verifiers in
different components to share spend state, i.e.\ to communicate, contradicting the
partition. Hence the tight achievable guarantee for an offline scheme is a
\emph{per-verifier} cap of $B$; a per-\emph{component} cap of $B$ additionally
requires (reachable) intra-component coordination; a \emph{global} cap of $B$ is
unachievable without cross-component consensus.
\end{proposition}

\begin{proof}[Justification]
The upper construction is immediate: the same signed token verifies identically
at each gateway, and offline meters are independent, so each admits up to $B$.
For the lower bound, suppose an offline scheme guaranteed total consumption
$< kB$. Consider two honest gateways $g_i\in C_i$, $g_j\in C_j$ that never
communicate. Each must decide admission from only the token and its local state;
by an indistinguishability argument, $g_i$'s decisions are independent of any
spend at $g_j$, so if $g_i$ would admit $B$ in isolation it admits $B$ here ---
and symmetrically $g_j$ --- for total $\geq 2B$ across just these two, and $kB$
across all $k$. Preventing this needs $g_i,g_j$ to share state, i.e.\ a
communication path, contradicting partition. This is precisely the offline
double-spending phenomenon of untraceable e-cash~\cite{chaum1990ecash}: an
offline bearer instrument cannot be prevented from double-spending, only
detected.
\end{proof}

To our knowledge this framing --- \emph{mesh egress quota is offline
double-spending} --- has not been made explicit, and it is clarifying, because it
transfers a 35-year-old resolution. Chaum--Fiat--Naor do not prevent
double-spending offline (they prove one cannot); they \emph{detect and identify}
the double-spender after the fact, and that threat deters. The same design fits
here:

\begin{quote}
\textbf{Achievable design.} (i)~Each gateway hard-caps a capability at $B$
locally (\emph{prevention within a verifier} --- what we implement, PR-level
\texttt{QuotaLedger}). (ii)~Gateways gossip signed spend-receipts
$(\text{cap id},\text{bytes},\text{gateway})$ as ordinary mesh traffic. When
partitions merge, a capability whose summed spend exceeds $B$ is
\emph{provably over-spent}, and --- if issuance binds the holder's identity --- the
over-spender is identified and revocable (\emph{detection and attribution across
components}). The system thus degrades gracefully from an impossible ``hard global
bound'' to ``hard per-verifier bound $+$ eventual global detection,'' with no
online coordination.
\end{quote}

We implement (i) and the revocation hook; the spend-gossip detector of (ii) is
specified but not yet built, and is the most concrete piece of future work the
paper produces.

\subsection{Sealed-sender aligns privacy with accountability}\label{sub:privacy}

Reputation systems normally fight Sybil \emph{badmouthing}: an adversary mints
identities to defame honest nodes by injecting false negative scores into shared
reputation state. The usual view is that privacy and accountability trade off ---
more hiding, less attributability. Sealed-sender delivery inverts this.

\begin{proposition}[Privacy--accountability alignment]\label{prop:privacy}
In a sealed-sender DTN where a delivery receipt is decryptable only by the
original sender, any custody-reputation signal is necessarily computed at the
source (\S\ref{par:custody}), from receipts that are unforgeable and sealed to it.
Therefore there is no reputation state shared between nodes. A Sybil adversary has
nothing to badmouth or inflate: it cannot write to another node's source-local
view, and forging a receipt requires the sender's key. The confidentiality
mechanism (hiding the sender) removes the shared-state attack surface that
reputation systems must otherwise defend, so here privacy and reputation-robustness
are complementary rather than opposed.
\end{proposition}

The cost is that a node learns only about \emph{its own} traffic's fate ---
reputation does not aggregate across the network. For the black-hole/grey-hole
threat this is sufficient (each source independently routes around peers that eat
\emph{its} bundles); a network-wide reputation would need a poisoning-resistant
aggregation layer we deliberately omit.

\subsection{The differential lens (a framing, not a finding)}\label{sub:diff}

Several bandwidth optimizations are the same idea --- send only the residual
$\Delta = Q - \mathbb{E}[Q\mid R]$ against a shared reference $R$ --- in two dual
regimes: a lossless-coding regime (set reconciliation transmits $O(d\log|U|)$ for
a size-$d$ difference~\cite{minsky2003setrec}; a manifest transmits only missing
CIDs; the held-bundle digest is re-sent only when its fingerprint moves) and an
estimation regime (mesh time transmits an offset against a stratum reference, not
an absolute clock). We record this because the single lens saved re-derivation and
a single validity check against the information-theoretic literature ---
\emph{but we are explicit that it is a design lens, not a result}: every
underlying bound is classical. We call it out precisely so it is not mistaken for
a contribution.

\subsection{What is off-the-shelf versus new}\label{sub:composition}

\begin{table}[t]\centering\small
\begin{tabular}{@{}p{3.4cm}p{4.4cm}p{4.8cm}@{}}
\toprule
\textbf{Component} & \textbf{Basis (off-the-shelf)} & \textbf{What, if anything, is new}\\
\midrule
Identity & self-certifying keys & nothing (standard)\\
Transfer integrity/dedup & content addressing, Merkle~\cite{benet2014ipfs,merkle1987digsig} & nothing (standard)\\
Multi-source fetch & BitTorrent swarm~\cite{cohen2003bittorrent} & nothing of note (DTN adaptation)\\
Routing & spray-and-wait~\cite{spyropoulos2005sprayandwait} & nothing (standard)\\
Message crypto & sealed box, prekeys, X3DH~\cite{marlinspike2016x3dh} & nothing; honest FS caveat\\
Custody reputation & DTN reputation & \emph{source-locality forced by sealed-sender} $\Rightarrow$ poisoning-immunity (Prop.~\ref{prop:privacy})\\
Egress authorization & macaroons/Biscuit~\cite{birgisson2014macaroons,biscuit2021} & \emph{open/gated asymmetry} (Obs.~\ref{obs:asym}); \emph{quota = double-spending} + detection design (Prop.~\ref{prop:doublespend})\\
\bottomrule
\end{tabular}
\caption{An honest accounting. The system is a composition; the contribution is
the three right-column items in \S\ref{sec:obs}, all conceptual.}
\label{tab:composition}
\end{table}

Table~\ref{tab:composition} is the paper in one view: almost everything is
standard, and the value is in three observations at the seams, not in any part.

% ===========================================================================
\section{Security Analysis}\label{sec:security}
% ===========================================================================

We analyze against a taxonomy from STRIDE~\cite{shostack2014threatmodeling} and
the privacy dimensions of LINDDUN~\cite{deng2011linddun}, plus overlay/DTN-specific
attacks (Table~\ref{tab:threats}). The cautionary reference for this class of
system is the break of Bridgefy~\cite{albrecht2021bridgefy}, a deployed BLE mesh
messenger with neither authenticity nor confidentiality; its failures motivate our
sealed-sender design and refusal to trust relays. We treat an independent
cryptographic and protocol audit as a launch gate.

\begin{table}[t]\centering\small
\begin{tabular}{@{}p{3.0cm}p{6.0cm}p{4.1cm}@{}}
\toprule
\textbf{Threat} & \textbf{Mechanism} & \textbf{Residual limitation}\\
\midrule
Sybil~\cite{douceur2002sybil} & no shared reputation to inflate (Prop.~\ref{prop:privacy}); source-local attribution & DHT population still Sybil-able; no proof-of-personhood\\
Eclipse~\cite{singh2006eclipse,heilman2015eclipse} & Kademlia incumbent retention $+$ per-endpoint bucket cap & id-space Sybil possible; no PoW node ids\\
Wormhole~\cite{hu2003leashes} & per-CID content integrity; onion hides graph & no packet leashes\\
Black/grey-hole & source-attributed reputation: consecutive-miss $+$ lifetime-ratio demotion & heuristic; adaptive adversary near the ratio threshold\\
Spoofing/tampering & Ed25519 signature inside the sealed envelope; per-CID integrity & key seizure defeats it\\
DoS/spam & priority-scaled postage~\cite{back2002hashcash}; bounded buffers & sustained regional flooding\\
Egress over-use & per-verifier hard cap (Prop.~\ref{prop:doublespend}) & global bound impossible offline; detector unbuilt\\
Info disclosure & sealed-sender E2E; onion & gateway sees egress plaintext (by design)\\
Metadata/traffic analysis & sealed sender, onion, unlinkable bridge subkeys & \emph{not} resistant to a global adversary\\
Egress SSRF & scheme/host guards $+$ resolved-IP recheck~\cite{owasp_ssrf} & depends on the fetcher honoring the contract\\
Token forgery/escalation & signed, chained, offline-verifiable capabilities & format needs independent review\\
\bottomrule
\end{tabular}
\caption{Threats, mechanisms, and honestly-stated residuals. Implementation
hardening (bounded CBOR decode, decompression-bomb limits, zeroization, domain
separation) omitted for space.}
\label{tab:threats}
\end{table}

% ===========================================================================
\section{Evaluation}\label{sec:eval}
% ===========================================================================

Lifeline ships an acceptance simulation held to $\geq 95\%$ delivery across
partition, churn, and loss, plus $\sim$210 unit/integration tests. That
establishes functional correctness and a delivery floor, not comparative
performance. For the latter we ship a reproducible harness (\texttt{sim::bench})
that runs the \emph{same} scenario under three forwarding strategies and reports
the DTN standard metric set --- delivery ratio, latency, overhead (relay) ratio,
hop count, buffer occupancy~\cite{spyropoulos2005sprayandwait,vahdat2000epidemic}.

Table~\ref{tab:results} is a three-cluster partition bridged by one mule.
Direct delivery (no relay) cannot cross the partition; epidemic~\cite{vahdat2000epidemic}
floods everything at the highest cost; binary spray-and-wait matches epidemic's
delivery and latency at $2.4\times$ lower overhead \emph{and} buffer.
Table~\ref{tab:rwp} is the canonical Random Waypoint
setting~\cite{broch1998comparison} (24 nodes) --- a connected regime where even
direct delivery reaches everyone, but $5.5\times$ slower than spray-and-wait,
while epidemic is marginally fastest at $4\times$ the buffer and overhead.
Spray-and-wait is the knee of the curve in both.

\begin{table}[t]\centering\small
\begin{tabular}{@{}lrrrrr@{}}
\toprule
\textbf{Strategy} & \textbf{Delivery} & \textbf{Median lat.} & \textbf{Overhead} & \textbf{Mean hops} & \textbf{Peak buf.}\\
\midrule
Direct (no relay) & $0.0\%$ & --- & $0.00$ & --- & $24$\\
Spray-and-wait \textbf{(default)} & $100.0\%$ & $90$\,s & $\mathbf{10.0}$ & $2.75$ & $\mathbf{240}$\\
Epidemic & $100.0\%$ & $90$\,s & $24.0$ & $2.75$ & $576$\\
\bottomrule
\end{tabular}
\caption{Mule-partition scenario (three isolated clusters, one mule; seed 42, 700
ticks). Overhead $=$ forwarded copies per delivered message; buffer $=$ peak total
bundles across all nodes. All delivered messages were offline-verified.}
\label{tab:results}
\end{table}

\begin{table}[t]\centering\small
\begin{tabular}{@{}lrrrrr@{}}
\toprule
\textbf{Strategy} & \textbf{Delivery} & \textbf{Median lat.} & \textbf{Overhead} & \textbf{Mean hops} & \textbf{Peak buf.}\\
\midrule
Direct (no relay) & $100.0\%$ & $220$\,s & $2.0$ & $1.00$ & $\mathbf{48}$\\
Spray-and-wait \textbf{(default)} & $100.0\%$ & $40$\,s & $\mathbf{11.5}$ & $2.21$ & $276$\\
Epidemic & $100.0\%$ & $\mathbf{30}$\,s & $46.0$ & $4.04$ & $1104$\\
\bottomrule
\end{tabular}
\caption{Random Waypoint (24 nodes, seed 42, 700 ticks): a connected regime.}
\label{tab:rwp}
\end{table}

The harness also replays an explicit contact trace (a \texttt{TraceMobility}
driver), so real traces --- CRAWDAD \texttt{cambridge/haggle}~\cite{crawdad_haggle}
and Reality Mining~\cite{eagle2006realitymining}, not vendored here for licensing
--- can be fed in verbatim, and cross-checked against the reference ONE
simulator~\cite{keranen2009one} and the Working Day
model~\cite{ekman2008workingday}. That external-trace study, and building the
spend-gossip detector of \S\ref{sub:doublespend}, are the principal remaining
empirical work.

% ===========================================================================
\section{Related Work}\label{sec:related}
% ===========================================================================

Lifeline sits at the confluence of DTN routing~\cite{rfc4838,vahdat2000epidemic,
spyropoulos2005sprayandwait}, content-addressed storage~\cite{benet2014ipfs,
cohen2003bittorrent}, secure messaging~\cite{marlinspike2016x3dh,
perrin2016doubleratchet}, overlays~\cite{maymounkov2002kademlia}, and
capability authorization~\cite{birgisson2014macaroons,biscuit2021}. It differs
from deployed mesh messengers in refusing to trust relays (cf.\ the Bridgefy
break~\cite{albrecht2021bridgefy}) and from Meshtastic~\cite{meshtastic} in
sealed-sender confidentiality plus a routing/custody layer. Our egress-quota
result imports offline e-cash~\cite{chaum1990ecash}: the connection between
\emph{bearer-capability metering} and \emph{offline double-spending} is, to our
knowledge, not previously drawn for the mesh setting, and it is what lets us state
what is and is not achievable rather than merely listing a TODO. Anonymity goals
are weaker than mixnets~\cite{vandenhooff2015vuvuzela,piotrowska2017loopix}, a
deliberate trade for latency and offline operation.

% ===========================================================================
\section{Limitations and Future Work}\label{sec:future}
% ===========================================================================

\begin{itemize}[leftmargin=1.6em,itemsep=1pt]
  \item \textbf{Independent review} --- no external cryptographic/protocol audit;
  the capability format and the forward-secrecy claim need one. This gates any
  real deployment.
  \item \textbf{The spend-gossip detector} (\S\ref{sub:doublespend}) is specified
  but unbuilt; it is the concrete realization of the paper's main observation.
  \item \textbf{External-trace evaluation} --- the harness is ready; real
  CRAWDAD/Reality-Mining runs and an ONE cross-check remain.
  \item \textbf{Metadata protection} --- traffic-analysis resistance is limited;
  cover-traffic/mixing~\cite{piotrowska2017loopix} is future work.
  \item \textbf{Overlay hardening} --- the DHT caps contacts per endpoint per
  bucket and custody attribution catches grey-holes, but proof-of-work node ids
  and adaptive-adversary resistance remain open.
\end{itemize}

% ===========================================================================
\section{Conclusion}
% ===========================================================================

Project Lifeline shows a messenger can keep its identity, delivery guarantee,
confidentiality, and \emph{authorization} working when every server is gone --- by
moving each into the endpoints and the data. We have tried to be honest that the
machinery is a composition of known primitives. The value is in three things that
composition exposed: that an open relay commons beneath a gated egress is the
right shape; that metering that gate offline \emph{is} double-spending, so the
achievable guarantee is per-component with detection-on-merge, not a global bound;
and that sealed-sender delivery makes privacy and reputation-robustness
complementary. The first of these turns an engineering dead-end into a principled
design; it is the one we would most like others to build on.

\small
\bibliographystyle{plain}
\bibliography{lifeline}

\end{document}
